\chapter{The Habitable Zone and the HZIED Test}
\label{chap:exoplanet_characterization}


We now aim to setup the physics and observational vocabulary for this work. \Cref{sec:habitable_zone} defines the classical
habitable zone and locates its inner edge at the runaway greenhouse
threshold.  \Cref{sec:transit_method} covers the transit method and the
precision it delivers on radii of small planets.
\Cref{sec:small_planet_demographics} sketches the small-planet
population.  \Cref{sec:hzied_framework} closes by defining
the \HZIED test itself, the specific hypothesis that
\textcite{Schlecker2024} proposed and this thesis extends.

% (write section body here)

\section{Habitable Zone}
\label{sec:habitable_zone}


With the habitable zone we refer to the circumstellar region where a earth-like rocky planet (holding a CO$_2$/H$_2$O/N$_2$ atmosphere regulated by a functioning
carbonate--silicate cycle) can sustain an amount of liquid surface water \parencite{Kasting1993,Kopparapu2013} \openquestion{Is this definition correct?}.
It is bounded by the inner edge where a water-saturated atmosphere can no longer radiate the incoming stellar flux (instellation) back to space, and the runaway greenhouse effect sets in.
The outer side is bounded by the classical maxium-greenhouse limit of \textcite{Kasting1993} and \textcite{Kopparapu2013}:

CO$_2$ condenses out of the atmosphere and the greenhouse effect is no longer sufficient to keep the surface temperature above the freezing point of water \parencite{Kopparapu2013}.
Both boundaries, functions of stellar type and atmospheric composition, are  theoretical rather then observed. They emerge from 1D radiative-convective \openquestion{what is this exactly; worth mentioned here?} climate models assuming earth-like compositions.


Unlike the inner edge, the outer edge is not a sharp threshold but rather
a gradual transition from habitable to uninhabitable conditions.  It can be
extended by non-condensing greenhouse gases: primordial H$_2$--He mixtures
accreted from protoplanetary disks can prevent the surface from freezing
\parencite{PierrehumbertGaidos2011}, while continuous volcanic outgassing of
H$_2$ pushes the classical outer edge to ${\sim}\,2.4\,\mathrm{AU}$ in the solar
system \parencite{RamirezKaltenegger2017}.  Volcanic CH$_4$ can similarly widen
the outer HZ by ${\sim}\,8\,\%$ for a solar-analog host, though it cools rather
than warms atmospheres around cooler M-dwarfs (see \textcite{Ramirez2018} for a
review).


Most importantly for this thesis, the outer transition leaves no clean signature
on the transit-radius plane.  A cold rocky planet near the maximum-greenhouse
limit has essentially the same bulk radius as its temperate counterpart: no
atmospheric phase transition alters its envelope mass or bulk structure in a
detectable way.
he runaway greenhouse, by contrast, inflates transit radii by
tens of percent (\Cref{sec:hzied_framework}), producing a step function in the
population-averaged radius distribution.  The \HZIED test therefore focuses
exclusively on the inner edge, and the outer boundary appears in the rest of
this thesis only as context.

\section{The Habitable Zone Inner-Edge Discontinuity}
\label{sec:hzied_framework}

% (write section body here)

The HZIED is the statistical signature of the runaway greenhouse effect on the transit-radius of a population of small planets.
This claim could be divided into three lines: (i) a small rocky planet with an amount of surface water crossing the runaway threshold $S_\mathrm{thresh}$ inflates its transit-observable radius by tens of percent ; (ii) at a fixed water inventory, the change of radius reveals a step function in instellation S and (iii) this step function survives averaging over an entire population of rocky planets depsite noisy detection. The detection of the HZIED would therefore be a strong evidence for the runaway greenhouse transition on exoplanets.




Three physical ingredients had to be assembled before the idea became
a testable prediction.  \textcite{Turbet2019} first pointed out that
a steam-dominated atmosphere on a runaway planet is transit-observable,
with a radius inflation of tens of percent for Earth-mass planets.
\textcite{Turbet2020} then computed a full mass--radius grid for
 rocky planets more irradiated than the runaway greenhouse
limit, tabulating the inflated transit radius as a function of planet
mass and water mass fraction $x_{\mathrm{H_2O}}$.
\textcite{DornLichtenberg2021} added the constraint that a substantial
fraction of the outgassed water dissolves back into the silicate melt
during the magma-ocean phase and does not contribute to the
atmospheric expansion, reducing the effective radius inflation by up to a
factor of two at $x_{\mathrm{H_2O}} \sim 10^{-2}$.  Together these
three form the ``Turbet+2020 with Dorn--Lichtenberg (DL) correction'',
forming the original \HZIED framework.





% (write section body here)
\section{Transit Method}
\label{sec:transit_method}

% (write section body here)

\section{Small-Planet Demographics}
\label{sec:small_planet_demographics}

% (write section body here)

